\chapter{Fragmentos de código relevantes}
\label{app:codigo}

Este apéndice selecciona principalmente fragmentos de \path{experiments/proyecto3_simulation.py}. El ejemplo de bootstrap reproduce un control post hoc aplicado a los CSV y se identifica como tal; no pertenece al generador autoritativo. El objetivo es mostrar cómo las decisiones metodológicas se traducen a código. Los recortes no reemplazan los archivos ejecutables ni sus validaciones.

\section{Parámetros y semillas}

\begin{lstlisting}[language=Python, caption={Contrato de parámetros de una corrida.}]
@dataclass
class SimParams:
    distance_km: float
    attenuation_db_km: float = DEFAULT_ALPHA_DB_KM
    detector_efficiency: float = 0.1
    dark_count_hz: float = 100.0
    mean_photon_num: float = 0.1
    frequency_hz: float = DEFAULT_FREQUENCY_HZ
    visibility: float = 0.98
    key_length: int = DEFAULT_KEY_LENGTH
    num_keys: int = DEFAULT_NUM_KEYS
    runtime_ps: float = DEFAULT_RUNTIME_PS
    alice_seed: int = 0
    bob_seed: int = 1
    count_rate_hz: float = 50e6
    time_resolution_ps: int = 10
    detection_window_ps: int = DEFAULT_DETECTION_WINDOW_PS
    classical_extra_delay_ps: int = DEFAULT_CLASSICAL_EXTRA_DELAY_PS
\end{lstlisting}

\texttt{SimParams} separa la configuración del motor de simulación. Cada barrido crea una instancia y modifica sólo los campos definidos por su diseño. La ventana de detección aparece en la misma estructura, pero se usa en la referencia analítica de fondo; no configura una compuerta física de SeQUeNCe. Esta diferencia evita atribuir al receptor una función que no implementa.

\begin{lstlisting}[language=Python, caption={Semilla común para las fuentes aleatorias del modelo.}]
def _seed_simulation(alice_seed: int, bob_seed: int) -> int:
    simulation_seed = (
        (alice_seed * 1_000_003) ^ bob_seed
    ) & 0xFFFFFFFF
    np.random.seed(simulation_seed)
    random.seed(simulation_seed)
    Timeline.seed(simulation_seed)
    return simulation_seed
\end{lstlisting}

La corrida usa varias fuentes de aleatoriedad. Sembrar sólo NumPy no alcanza: también se fijan el módulo estándar y la línea temporal. Alice y Bob reciben, además, sus semillas de nodo. El flujo asigna pares no superpuestos por punto y repetición. Así puede reconstruirse una observación sin hacer idénticas todas las corridas.

\section{Enlace y parche de visibilidad}

\begin{lstlisting}[language=Python, caption={Creación abreviada del enlace y configuración del receptor.}]
tl = Timeline(p.runtime_ps)
qc0 = QuantumChannel(
    "qc0", tl, distance=distance_m,
    attenuation=attenuation_db_m,
    frequency=p.frequency_hz,
)
cc0 = ClassicalChannel("cc0", tl, distance=distance_m)
cc0.delay += p.classical_extra_delay_ps

alice = QKDNode("alice", tl, encoding=time_bin, stack_size=1)
bob = QKDNode("bob", tl, encoding=time_bin, stack_size=1)

phase_error = max(0.0, min(1.0, (1.0 - p.visibility) / 2.0))
qsd = bob.components["bob.qsdetector"]
qsd.update_interferometer_params("phase_error", phase_error)
\end{lstlisting}

La distancia y atenuación se aplican al canal cuántico, mientras el retardo adicional pertenece al canal clásico. El experimento de detector fija este último en cero para reducir su influencia sobre la tasa. La visibilidad se transforma en error de fase y llega al interferómetro.

La versión base de SeQUeNCe 0.8.5 no aplicaba correctamente ese error a \texttt{FreeQuantumState}. El repositorio corrige \path{sequence/components/interferometer.py}. Sin el parche, barrer \(V\) no produciría la respuesta documentada; por eso su hash forma parte del manifiesto.

\section{Tasa agregada y referencia}

\begin{lstlisting}[language=Python, caption={Cálculo de la tasa tamizada de una repetición.}]
bb_a = alice.protocol_stack[0]
errs = list(getattr(bb_a, "error_rates", []) or [])
n_keys = len(errs)
total_sifted_bits = n_keys * p.key_length
total_errors = int(round(sum(errs) * p.key_length))
mean_qber = (
    total_errors / total_sifted_bits
    if total_sifted_bits else float("nan")
)

elapsed_key_ps = float(getattr(bb_a, "last_key_time", 0.0))
elapsed_key_s = elapsed_key_ps * 1e-12
aggregate_throughput = (
    total_sifted_bits / elapsed_key_s
    if elapsed_key_s > 0 else 0.0
)
\end{lstlisting}

SeQUeNCe expone una tasa de error por clave completada en \texttt{error\_rates}; no expone aquí las claves tamizadas como atributos públicos. El script reconstruye el conteo total de errores multiplicando la suma de esas tasas por la longitud de clave y redondeando al entero más cercano. Luego calcula QBER con los conteos agregados, no promediando porcentajes por clave. La tasa divide todos los bits tamizados por el tiempo simulado hasta la última clave. Promediar tasas instantáneas asignaría el mismo peso a intervalos de distinta duración y podría sesgar la salida.

\begin{lstlisting}[language=Python, caption={Fondo, ganancia y referencia de clics tamizados.}]
def background_yield(dark_count_hz, detection_window_ps,
                     detectors=3):
    exposure_s = detectors * detection_window_ps * 1e-12
    return 1.0 - math.exp(-dark_count_hz * exposure_s)

def wcs_gain(mu, eta_channel, eta_det, y0):
    no_signal_click = math.exp(-mu * eta_channel * eta_det)
    return 1.0 - (1.0 - y0) * no_signal_click

y0 = background_yield(p.dark_count_hz, p.detection_window_ps)
q_mu = wcs_gain(p.mean_photon_num, eta_ch,
                p.detector_efficiency, y0)
sifted_rate_reference = 0.5 * min(
    p.frequency_hz * q_mu,
    3 * p.count_rate_hz,
)
\end{lstlisting}

\texttt{background\_yield} convierte tasa oscura y ventana supuesta en probabilidad por pulso para tres detectores. \texttt{wcs\_gain} combina ese fondo con la probabilidad de no recibir señal. El factor \(1/2\) representa tamizado por bases y el mínimo impone disponibilidad de pulsos o límite agregado de conteo. La referencia es un control medio; no se programa como tope de cada corrida.

\section{Indicador ideal}

\begin{lstlisting}[language=Python, caption={Indicador de posprocesamiento monofotónico ideal.}]
def ideal_postprocessing_rate(qber, r_sifted_bps,
                              f_ec=1.16):
    if (not math.isfinite(qber)
            or qber < 0 or qber >= 0.5):
        return 0.0
    h = _binary_entropy(qber)
    factor = max(0.0, 1.0 - f_ec * h - h)
    return r_sifted_bps * factor
\end{lstlisting}

La función no fija un corte manual en 11\%. El factor se anula por su propia entropía cerca de 9,81\% cuando \(f_\mathrm{EC}=1{,}16\). Su nombre evita una afirmación que el cálculo no sostiene: para una fuente coherente débil sin estimación monofotónica, el resultado no es una tasa secreta.

\section{Intervalos y contrastes}

\begin{lstlisting}[language=Python, caption={Intervalo de Wilson para QBER agrupada.}]
def _wilson_qber_ci(runs, confidence=0.95):
    errors = sum(run["total_errors"] for run in runs)
    bits = sum(run["total_sifted_bits"] for run in runs)
    p = errors / bits
    z = stats.norm.ppf((1.0 + confidence) / 2.0)
    denom = 1.0 + z * z / bits
    center = (p + z * z / (2.0 * bits)) / denom
    half = (
        z * math.sqrt(
            p * (1.0 - p) / bits
            + z * z / (4.0 * bits * bits)
        ) / denom
    )
    return p, max(0.0, center-half), min(1.0, center+half)
\end{lstlisting}

QBER es una proporción y puede estar cerca de cero. Wilson conserva mejor comportamiento que la aproximación normal simple en ese régimen, siempre que los bits agrupados puedan tratarse como intercambiables. No captura dependencia dentro de una clave ni heterogeneidad entre repeticiones. Las tasas usan intervalos t porque cada repetición produce una observación continua; las diferencias de detector usan Welch para no imponer varianzas iguales.

Como control posterior, las medias señuelo se remuestrean a nivel de repetición:

\begin{lstlisting}[language=Python, caption={Bootstrap percentil usado como análisis de sensibilidad.}]
rng = np.random.default_rng(20260901)
indices = rng.integers(
    0, rates.size, size=(100_000, rates.size)
)
bootstrap_means = rates[indices].mean(axis=1)
ci_bootstrap = np.quantile(
    bootstrap_means, [0.025, 0.975]
)
\end{lstlisting}

Este cálculo usa las 30 tasas ya exportadas y no vuelve a ejecutar SeQUeNCe. El remuestreo conserva como unidad la repetición completa, por lo que tolera asimetría y valores nulos mejor que una justificación puramente normal \cite{efron1979bootstrap}. Para visibilidad se usa además HC3, que pondera cada residuo por \((1-h_{ii})^{-2}\) y reduce la dependencia del intervalo respecto de homoscedasticidad \cite{mackinnonWhite1985}. Ninguno de estos controles convierte los resultados en cotas de clave finita.

\section{Cota señuelo y caso conservador}

\begin{lstlisting}[language=Python, caption={Cota de error monofotónico con reemplazo conservador.}]
def decoy_e1_upper(e_nu, q_nu, e0, y0, y1, nu):
    if y1 <= 0 or nu <= 0:
        return 0.5
    numerator = e_nu * q_nu * math.exp(nu) - e0 * y0
    if numerator <= 0:
        return 0.5
    return min(0.5, numerator / (y1 * nu))
\end{lstlisting}

Un numerador no positivo aparece cuando la combinación de error simulado, ganancia analítica y muestra finita no sostiene una cota positiva. Truncarlo a cero supondría ausencia de error y aumentaría la tasa. El código adopta \(1/2\), que anula la contribución secreta de esos eventos. El flujo cuenta cada activación para que la media no oculte esta decisión.

\section{Validaciones y manifiesto}

\begin{lstlisting}[language=Python, caption={Ejemplos de invariantes verificados antes de exportar.}]
assert 0.0 <= qber <= 0.5
assert qber_low <= qber <= qber_high
assert 0.0 <= ideal_rate <= sifted_rate
assert 0.0 <= q_nu <= q_mu <= 1.0
assert detector_sum == total_detector_clicks
for detector in qsd.detectors:
    for name, expected in effective_detector_params.items():
        if getattr(detector, name) != expected:
            raise RuntimeError(f"{name} was not applied")
\end{lstlisting}

Las aserciones detienen la generación si una salida viola las definiciones implementadas. No comprueban todavía que cada ejecución complete exactamente el número de claves solicitado. Después, el script calcula SHA-256 del flujo, el parche, las dependencias, cada CSV y las macros numéricas. Los hashes no prueban corrección científica ni cubren las fuentes LaTeX, las figuras o el PDF; prueban identidad de la evidencia computacional registrada.
